% Section body only: the master supplies \section{The Supersonic Base-Drag Closure
% and Its Components} and \input's this file. Start at \subsection.
% Content: the supersonic base-drag closure and its component branches, each with
% governing equation, regime of validity, and source; plus the C^1-continuous
% transonic blend. We DESCRIBE the closure and note that alternative literature
% closures exist; we DO NOT swap them one at a time and report per-closure apogee
% deltas (that experiment is out of scope). Mined from the ChapmanKorstBaseDrag /
% BarrowmanDragCalculator code.
% CITATION RULE (CANONICAL_FACTS.md Sec. H): the empirical Cd_base = 0.064 + 0.186/M^2
% correlation is presented as EMPIRICAL, validated vs NACA TN 3393 and consistent with
% ESDU 77021. NO Devan-Ashwood. All numbers per _shared/CANONICAL_FACTS.md.
\label{sec:basedrag-models}

This section specifies the supersonic base-drag closure as implemented and its
component branches. Each branch contributes to the base-drag coefficient
$C_{d,\text{base}}$ referenced to the base area
$A_{\text{base}} = \pi R_{\text{base}}^{2}$, so that its contribution to the
vehicle axial-force coefficient is $C_{d,\text{base}}\,A_{\text{base}}/A_{\text{ref}}$.
For each branch we state its governing equation, its regime of validity, and the
published source against which it was assessed. Three of the six components are
geometry-independent functions of Mach number and boundary-layer state --- the
turbulent Chapman--Korst shear-layer model, the Chapman laminar correlation, and
the empirical supersonic correlation --- consolidated under the ESDU~77021 data
item; the remaining three are geometry modifiers applied multiplicatively to the
clean-body value: the ESDU~77021 boundary-layer-thickness correction, the
Viswanath boattail relief, and the finned-base augmentation. The branches are
stitched together below Mach~$1.5$ by a single $C^{1}$-continuous transonic
blend, described last. Table~\ref{tab:closures} collects the inventory. The
engineering literature offers \emph{alternative} closures for several of these
branches --- competing turbulent base-pressure correlations, different
boattail-relief models, and several finned-base treatments --- that disagree in
functional form and were calibrated on different geometry families; we note these
alternatives where relevant, but this study describes and benchmarks the closure
as implemented rather than substituting alternatives one at a time and reporting
per-closure apogee deltas, which is outside its scope.

\subsection{Empirical supersonic correlation}
\label{subsec:empirical}

The workhorse turbulent closure is the two-coefficient empirical correlation
%
\begin{equation}
C_{d,\text{base}} = 0.064 + \frac{0.186}{M^{2}}, \qquad M > 1.3 ,
\label{eq:empirical}
\end{equation}
%
applicable to a slender axisymmetric afterbody with a turbulent boundary layer at
separation. The two coefficients encode the two physical asymptotes of supersonic
base flow: the leading constant $0.064$ is the nonzero high-Mach floor toward
which the base-pressure coefficient asymptotes as the freestream entrains the
wake, and the $0.186/M^{2}$ term captures the decay of the base-pressure deficit
with dynamic pressure as Mach number rises. We present
Eq.~\eqref{eq:empirical} strictly as an empirical fit: it reproduces the
turbulent NACA~TN~3393~\cite{chapman1955} nonlifting-body base-pressure
measurements over $M = 2.73$--$4.48$ to a mean absolute percentage error of
$15.9\%$ (and the Hart free-flight subset to $4.0\%$), and its functional form is
consistent with the recommended supersonic base-pressure correlation
consolidated in ESDU~77021~\cite{esdu77021}. No further provenance is claimed for
the two coefficients beyond this benchmark agreement.

\subsection{Chapman laminar base drag}
\label{subsec:laminar}

For perfect-finish vehicles whose surface is smooth enough that boundary-layer
transition is delayed and the layer remains laminar to the base, the turbulent
form of Eq.~\eqref{eq:empirical} over-predicts the base suction. The Chapman
laminar correlation~\cite{chapman1950} applies instead:
%
\begin{equation}
C_{p,b,\text{lam}} = \frac{C_{\text{lam}}}{M^{2}\sqrt{\mathrm{Re}_{L}}} ,
\qquad C_{\text{lam}} = 1870 ,
\label{eq:laminar}
\end{equation}
%
with $\mathrm{Re}_{L}$ the body-length Reynolds number. The $\mathrm{Re}_{L}^{-1/2}$
scaling follows from laminar free-shear-layer mixing theory --- a laminar layer
carries less momentum, mixes less, and recompresses to a lower base pressure than
a turbulent layer --- and the $M^{-2}$ factor again arises from normalizing the
base-pressure deficit by freestream dynamic pressure. The constant
$C_{\text{lam}} = 1870$ is a geometric-mean fit to the four condensation-corrected
laminar TN~3393 points ($M = 2.73$--$4.48$, $\mathrm{Re}_{L} \approx 4$--$6 \times
10^{6}$), giving a mean absolute percentage error of $4.4\%$ against those data.
Applying the turbulent form of Eq.~\eqref{eq:empirical} to the same laminar data
yields a mean absolute percentage error of $44\%$: the boundary-layer state at
separation is the dominant discriminator between the two branches, which is
precisely why a base-drag benchmark cannot treat ``the'' base-drag closure
as a single object. Equation~\eqref{eq:laminar} is blended into the turbulent
result over $M = 1.3$--$2.5$ by the Hermite smoothstep $s(t) = 3t^{2} - 2t^{3}$,
weighted by the local laminar fraction so that a fully turbulent body never sees
the laminar branch; a clamp at the near-vacuum limit $2/(\gamma M^{2})$ prevents
unphysical growth at low Reynolds number.

\subsection{Chapman--Korst turbulent shear-layer model}
\label{subsec:chapmankorst}

The Chapman--Korst dead-air recompression theory~\cite{chapman1955} provides a
turbulent base-pressure estimate from boundary-layer edge conditions, used here
as the supersonic-side closure that the empirical correlation blends \emph{into}.
The supersonic shear layer separating at the base corner reattaches on the wake
axis at a recompression point; the balance between corner expansion and
shear-layer recompression fixes the base pressure. For a thin turbulent boundary
layer the clean-body coefficient follows the same family as
Eq.~\eqref{eq:empirical} with an added $M^{-4}$ term, calibrated to the turbulent
cylindrical-afterbody tabulation of ESDU~77021~\cite{esdu77021}:
%
\begin{equation}
C_{d,\text{base}}^{\,\text{thin}} = 0.060 + \frac{0.190}{M^{2}}
   + \frac{0.005}{M^{4}} .
\label{eq:chapmankorst}
\end{equation}
%
This reproduces the ESDU anchor values $C_{d,\text{base}} \approx 0.145$, $0.110$,
$0.085$, and $0.070$ at $M = 1.5$, $2.0$, $3.0$, and $5.0$, respectively, which
agree closely with Eq.~\eqref{eq:empirical} for a thin layer --- as expected,
since both describe the same recompression physics. Equation~\eqref{eq:chapmankorst}
is the baseline to which the boundary-layer-thickness correction of
Section~\ref{subsec:esdu} is applied. It is blended into the empirical correlation
over $M = 1.2$--$1.4$ by the same smoothstep $s(t)$, so that pure
Eq.~\eqref{eq:empirical} governs below $M = 1.2$ and the Chapman--Korst form
governs above $M = 1.4$.

\subsection{ESDU~77021 boundary-layer-thickness correction}
\label{subsec:esdu}

ESDU~77021~\cite{esdu77021} consolidates a broad experimental corpus of base
pressure on bodies of revolution at supersonic and hypersonic Mach numbers
without base injection or combustion, and is the reference data item for the
turbulent closures above. Beyond fixing the thin-layer baseline of
Eq.~\eqref{eq:chapmankorst}, it supplies the dependence on boundary-layer
thickness: a thicker layer entrains more mass into the recirculating wake, raises
base pressure, and so \emph{reduces} base drag. The correction is applied as a
multiplicative factor on the thin-layer coefficient,
%
\begin{equation}
f_{\delta} = 1 - k(M)\sqrt{\theta/R_{\text{base}}} ,
\qquad k(M) = 0.8 + \frac{0.2}{M} ,
\label{eq:esdu-bl}
\end{equation}
%
where $\theta$ is the turbulent momentum thickness at the base, estimated from the
flat-plate correlation $\theta/x = 0.036\,\mathrm{Re}_{x}^{-1/5}$, and
$\theta/R_{\text{base}}$ is clamped to $[10^{-5}, 0.5]$ with $f_{\delta}$ clamped
to $[0.3, 1.0]$. For a typical rocket ($\theta/R_{\text{base}} \approx 0.01$) the
reduction is a few percent; for an unusually thick layer
($\theta/R_{\text{base}} \approx 0.1$) it reaches $20$--$30\%$. This correction is
the mechanism by which the slender, high-fineness airframes in the corpus carry a
distinct base-drag signature from short, blunt bodies, and it is the
literature-anchored counterpart to the corpus-frozen thick-boundary-layer
amplification disclosed in the companion study~\cite{yu2026jsr}.

\subsection{Viswanath boattail relief}
\label{subsec:viswanath}

When the afterbody tapers inward to a boattail, the converging flow turns through
a compression that raises base pressure and shrinks the base area, reducing base
drag below the blunt-base value. Following Viswanath~\cite{viswanath1996}, the
relief is applied as a multiplicative retention factor $\eta_{bt}$ on the
clean-body base drag, parameterized by the boattail half-angle
$\theta_{bt} = \arctan\!\big[(R_{\text{fore}} - R_{\text{aft}})/L_{bt}\big]$ and,
in the effective regime, by Mach number:
%
\begin{equation}
\eta_{bt} =
\begin{cases}
0.25 + 0.05\,\theta_{bt}, & \theta_{bt} < 6^{\circ},\\[1.0ex]
\big[\,0.55 + 0.04(\theta_{bt} - 6^{\circ})\,\big]\,\big[1 + 0.1\max(0, M-1)\big],
   & 6^{\circ} \le \theta_{bt} < 16^{\circ},\\[1.0ex]
0.95 - 0.05(\theta_{bt} - 16^{\circ}), & \theta_{bt} \ge 16^{\circ},
\end{cases}
\label{eq:viswanath}
\end{equation}
%
with $\eta_{bt}$ clamped to $[0, 1]$. For well-chosen boattail angles of
$6$--$16^{\circ}$ this yields a $15$--$40\%$ base-drag reduction, consistent with
the afterbody-drag-management survey of Viswanath~\cite{viswanath1996}. The factor
is applied to the wake \emph{downstream} of a boattail and is not also applied on
the boattail component itself, which would double-count the relief and can drive
the aft base drag to zero on steep imported geometries that still expose a finite
blunt exit area.

\subsection{Finned-base augmentation (externally calibrated, B-level)}
\label{subsec:finnedbase}

The dominant apogee mechanism (\S\ref{subsec:ablation}) is the augmentation of
base drag by the tail fins. Fins mounted at or just ahead of the base disrupt the
recompression of the body wake, deepening the base suction relative to the
clean-body value; aeroballistic-range data on the Basic Finner~\cite{dupuis1997}
and the afterbody-drag compendium of Hoerner~\cite{hoerner1965} consistently show
$40$--$60\%$ higher base drag for four-fin configurations than for the
corresponding smooth cylindrical afterbody over $M = 1.5$--$3$, which we use as an
external \emph{bound} on the plausible magnitude of the effect rather than as a
calibration target. The augmentation is applied as a multiplicative factor on the
post-relief base drag, written here in its full implemented form,
%
\begin{equation}
\Phi_{\text{fin}} = 1 + K_{\text{fin}}\,
   \frac{1 - e^{-n_{\text{fins}}/1.4}}{1 - e^{-4/1.4}}\,
   \operatorname{clamp}\!\Big(\frac{s_{\max}}{R_{\text{base}}},\,0.3,\,1\Big)\,
   f(M) ,
\qquad K_{\text{fin}} = 0.55 ,
\label{eq:finnedbase}
\end{equation}
%
where the fin count $n_{\text{fins}}$ enters through a saturating term normalized
to unity at the four-fin reference (wake disruption is not linear in fin count:
three corner-wake sectors already nearly divide the base shear layer, so a
three-fin body is not under-counted), $s_{\max}$ is the maximum fin span, and the
Mach ramp $f(M)$ rises from $0$ at $M = 0.2$ to $0.30$ at $M = 0.8$, then from
$0.30$ to $1$ over $M = 0.8$--$1.3$, holds through the low-supersonic band to
$M = 3.0$, then decays as $3/M$ above $M = 3.0$ as the fins become progressively
submerged in the body shock cone. At the four-fin reference (four rectangular
fins, $s_{\max}/R_{\text{base}} = 2.0$, where both geometry factors saturate to
unity) the scale constant $K_{\text{fin}} = 0.55$ gives a $\sim\!50\%$
augmentation at $M = 2$, inside the $40$--$60\%$ literature bound above; companion
scale constants of $1.00$ for rounded-trailing-edge fins and $1.35$ for fin-can
sleeve geometries cover the remaining corpus afterbody types.

We classify $K_{\text{fin}}$ as a \emph{B-level} constant in the specific sense
that it lacks validation against an isolated finned-base base-\emph{pressure}
dataset. It is \emph{not} corpus-frozen: the value $0.55$ is calibrated
externally against the Basic Finner aeroballistic-range data~\cite{dupuis1997}
so that it reproduces the $\sim\!50\%$ augmentation measured there at $M = 2$,
which also places it inside the $40$--$60\%$ over-clean-body band of the same
fixture and the Hoerner compendium~\cite{hoerner1965}. Its anchor is thus a
published external fixture, not flights drawn from the calibration corpus, so it
does not make the headline in-sample. The constant nonetheless does not yet carry
the A-level external validation of Eqs.~\eqref{eq:empirical} and
\eqref{eq:laminar}, which are checked against isolated base-pressure data, and
this B-to-A gap is the explicit limitation of the augmentation closure. The
distinction is material precisely because the mechanism ablation of
\S\ref{subsec:ablation} ranks this term as the dominant driver of integrated
apogee error, shifting mean corpus apogee error by $8.10$ percentage points when
removed --- about $54$ times the $0.15$-point effect of the
shock-geometry pre-pass. The two genuinely corpus-frozen constants in the
base-drag treatment belong to \emph{different} closures --- the
thick-boundary-layer base-drag amplifier ($\textsc{thick\_bl\_k}$) and the
slender-body pressure-drag term ($\textsc{slender\_body\_k}$) --- and it is those
two constants, not $K_{\text{fin}}$, that the decontaminated prospective holdout
of \S\ref{subsec:holdout} quarantines and shows to generalize (blind holdout MAE
$3.95\%$ below development $5.47\%$). The defense of $K_{\text{fin}}$ is instead
its external Basic Finner anchor and the component base-pressure benchmarks.
Elevating $K_{\text{fin}}$ from this total-drag anchor to A-level validation
against an isolated finned-base base-pressure dataset is identified in the
Conclusions as the single highest-value evidence gap raised by this work.

\subsection{Transonic blend and $C^{1}$ continuity}
\label{subsec:transonic}

Below $M = 1.5$ the supersonic closures above are not valid, and the base drag
instead rises to a transonic peak near $M \approx 1.05$, where a flat-based body
reaches $C_{d,\text{base}} \approx 0.25$. The transonic band $M = 0.85$--$1.5$ is
represented by a degree-4 polynomial anchored at three points: the peak value
$C_{d,\text{base}} = 0.25$ at $M = 1.05$, the free-flight plateau
$C_{d,\text{base}} = 0.230$ at $M = 1.30$ (just below the Hart
NACA~RM~L52E06~\cite{hart1952basedrag} finless free-flight reading of
$0.250 \pm 0.013$, within its digitization uncertainty), and a $C^{1}$ handoff to
the supersonic correlation at $M = 1.5$.
The polynomial is constructed to be continuous in both value and first derivative
at the regime boundaries, and the laminar, Chapman--Korst, and boattail blends of
the preceding subsections all use the same Hermite smoothstep
$s(t) = 3t^{2} - 2t^{3}$, which has zero derivative at its endpoints. This global
$C^{1}$ continuity is not cosmetic: a discontinuity in $C_{d,\text{base}}$ or its
slope near Mach~1 drives the fourth-order Runge--Kutta trajectory integrator to
oscillate whenever the vehicle repeatedly crosses the sonic line during boost and
coast, and the resulting integration noise would contaminate the very apogee
sensitivity this paper sets out to measure. The same continuity discipline is
applied across every Mach-regime transition in the underlying simulator, as
detailed in the companion paper~\cite{yu2026jsr}.

\begin{table*}[t]
\centering
\caption{Components of the supersonic base-drag closure as implemented. Each is
referenced to base area; geometry modifiers (ESDU~77021 BL correction, Viswanath,
finned-base) are applied multiplicatively to the clean-body coefficient. Status A
denotes external component validation against an isolated base-pressure dataset; B
denotes a disclosed constant lacking such an isolated dataset (here, the
finned-base augmentation, calibrated externally against the Basic Finner
total-drag fixture). Alternative literature closures exist for several branches
but are not swapped in or scored at apogee in this study.}
\label{tab:closures}
\small
\setlength{\tabcolsep}{4pt}
\begin{tabularx}{\textwidth}{l l X X c}
\toprule
Component & Regime & Governing form & Source & Status \\
\midrule
Empirical supersonic & $M > 1.3$, turbulent BL & Eq.~\eqref{eq:empirical},
   $0.064 + 0.186/M^{2}$ & NACA~TN~3393~\cite{chapman1955}, ESDU~77021~\cite{esdu77021} & A \\
Chapman laminar & $M = 1.3$--$2.5$, laminar BL & Eq.~\eqref{eq:laminar},
   $C_{\text{lam}}/(M^{2}\sqrt{\mathrm{Re}_{L}})$ & Chapman~\cite{chapman1950}, NACA~TN~3393~\cite{chapman1955} & A \\
Chapman--Korst turbulent & $M = 1.2$--$5{+}$ & Eq.~\eqref{eq:chapmankorst},
   $0.060 + 0.190/M^{2} + 0.005/M^{4}$ & ESDU~77021~\cite{esdu77021} & A \\
ESDU~77021 BL correction & $M > 1.2$, thick BL & Eq.~\eqref{eq:esdu-bl},
   $1 - k(M)\sqrt{\theta/R}$ & ESDU~77021~\cite{esdu77021} & A \\
Viswanath boattail & Any $M$, boattail geometry & Eq.~\eqref{eq:viswanath},
   $\eta_{bt}(\theta_{bt}, M)$ & Viswanath~\cite{viswanath1996} & A \\
Finned-base augmentation & $M = 0.8$--$5{+}$, finned body & Eq.~\eqref{eq:finnedbase},
   $1 + K_{\text{fin}}\,g(n)\,h(s_{\max}/R)\,f(M)$ & Basic Finner~\cite{dupuis1997}, Hoerner~\cite{hoerner1965} (bound) & B \\
Transonic blend & $M = 0.85$--$1.5$ & Degree-4 polynomial, peak $0.25$ at $M{=}1.05$
   & Hart~RM~L52E06~\cite{hart1952basedrag} anchor & A \\
\bottomrule
\end{tabularx}
\end{table*}
